(1) 单粒子本征能量 \(\varepsilon_{n\uparrow}=\varepsilon_n-\mu B\)，\(\varepsilon_{n\downarrow}=\varepsilon_n+\mu B\)。Pauli原理禁止两个相同自旋占据同一 \(n\)；允许两个反向自旋占据同一空间轨道 \(n=1\)。大磁场时可能填充 \(n=1,2\) 的两个上自旋态。

(2) 低场填充 \(1\uparrow,1\downarrow\)，总能 \(E_{\rm low}=2\varepsilon_1\)；高场填充 \(1\uparrow,2\uparrow\)，能 \(E_{\rm high}=\varepsilon_1+\varepsilon_2-2\mu B\)。两者相等于 \(\boxed{B_c=(\varepsilon_2-\varepsilon_1)/(2\mu)=3\pi^2\hbar^2/(4m\mu L^2)}\)。在临界场其他候选填充不更低。

(3) \(M=-\partial_BE_G\)，因此 \(M=0\) 对 \(B<B_c\)，\(M=2\mu\) 对 \(B>B_c\)。这是一阶零温基态能级交叉造成的理想磁化阶跃。\(B=B_c\) 时两种填充具有相同最低能量，磁化需指定具体简并基态或微小扰动选择。

(4) 固定外磁场 \(B\)，因 \(\varepsilon_n\propto L^{-2}\)，低场壁压力 \(P_{\rm low}=4\varepsilon_1/L\)；高场 \(P_{\rm high}=2(\varepsilon_1+\varepsilon_2)/L\)，Zeeman常数项不贡献机械压强。极化后迫使一粒子进入 \(n=2\) 较高动量态，压强增加。\(B_c\propto L^{-2}\)，缩小阱长会推高完全自旋极化门槛。
